% ============================================================
% Gesamtskript VWL --- Übungsaufgaben Lektion 3
% Konjunktur, Wirtschaftspolitik & Außenwirtschaft (Rahmenplan DIHK 1.1)
% 8 Aufgaben, 100 Punkte --- Lösungen in loesungen-l3.tex
% ============================================================

\subsection*{Aufgabe~3.1 --- Konjunkturlage richtig deuten (11 Punkte)}
\addcontentsline{toc}{subsection}{Aufgabe~3.1 --- Konjunkturlage richtig deuten}

Sie sind Mitarbeiter/-in in der Unternehmensplanung der \emph{Kranich Fördertechnik GmbH},
einem mittelständischen Hersteller von Förderbändern und Hebeanlagen. Die Geschäftsführung
will wissen, wie sie die Produktions- und Personalplanung für das kommende Jahr ausrichten
soll. Ihnen liegen folgende aktuelle Wirtschaftsdaten für Deutschland vor:

\begin{enumerate}[label=(\arabic*)]
  \item Der ifo-Geschäftsklimaindex für das Verarbeitende Gewerbe ist zum dritten Mal in
    Folge gestiegen.
  \item Die Auftragseingänge in der Industrie liegen 3,5~\% über dem Vorquartal.
  \item Die Kapazitätsauslastung der Industrie beträgt 76~\% (langjähriger Durchschnitt:
    83~\%).
  \item Die Arbeitslosenquote ist im letzten Monat erneut leicht gestiegen.
  \item Die Zahl der Unternehmensinsolvenzen liegt auf dem höchsten Stand seit mehreren
    Jahren.
\end{enumerate}

\begin{enumerate}[label=\alph*)]
  \item \textbf{Ordnen} Sie die fünf Daten jeweils den Früh-, Präsenz- oder
    Spätindikatoren \textbf{zu}. \hfill (5~P.)
  \item \textbf{Bestimmen} Sie anhand der Daten die aktuelle Konjunkturphase und
    \textbf{begründen} Sie Ihre Einschätzung. \hfill (3~P.)
  \item \textbf{Erläutern} Sie zwei Schlussfolgerungen, die die Kranich Fördertechnik GmbH
    aus Ihrer Einschätzung für ihre Planung ziehen sollte. \hfill (3~P.)
\end{enumerate}

\subsection*{Aufgabe~3.2 --- Inflation, Realzins und Lohnabschluss (14 Punkte)}
\addcontentsline{toc}{subsection}{Aufgabe~3.2 --- Inflation, Realzins und Lohnabschluss}

Sie sind Mitarbeiter/-in im Rechnungswesen der \emph{Holtkamp Logistik GmbH}. Für eine
Präsentation in der Geschäftsleitung sollen Sie die Auswirkungen der Preisentwicklung
aufbereiten. Der Verbraucherpreisindex (VPI, Basisjahr 2020~$=100$) lag im September des
Vorjahres bei 124,0 Punkten und im September dieses Jahres bei 127,1 Punkten
(angenommene Werte).

\begin{enumerate}[label=\alph*)]
  \item \textbf{Berechnen} Sie die Inflationsrate für diesen Zeitraum. \hfill (3~P.)
  \item Die Holtkamp Logistik GmbH hat vor einem Jahr 250.000~€ als Festgeld für ein Jahr
    zu einem Zinssatz von 3,4~\% p.\,a. angelegt. \textbf{Berechnen} Sie den Zinsertrag in
    Euro sowie (näherungsweise) den Realzins, und \textbf{beurteilen} Sie das Ergebnis
    kurz. Gehen Sie davon aus, dass die Inflationsrate aus a) für die gesamte
    Anlagedauer gilt; Steuern bleiben unberücksichtigt. \hfill (4~P.)
  \item Das monatliche Bruttogehalt eines Disponenten wurde im selben Zeitraum von
    3.200~€ auf 3.264~€ angehoben. \textbf{Berechnen} Sie die nominale Gehaltssteigerung
    in Prozent und \textbf{ermitteln} Sie näherungsweise die Veränderung des Reallohns.
    \textbf{Interpretieren} Sie das Ergebnis. \hfill (3~P.)
  \item In der anstehenden Tarifrunde fordert die Gewerkschaft eine Lohnerhöhung von
    5,5~\%. Die Arbeitsproduktivität ist im letzten Jahr um 1,3~\% gestiegen.
    \textbf{Berechnen} Sie den kostenniveauneutralen Verteilungsspielraum (Zielinflation
    der EZB) und \textbf{beurteilen} Sie die Forderung aus Sicht der Arbeitgeber und aus
    Sicht der Gewerkschaft. \hfill (4~P.)
\end{enumerate}

\subsection*{Aufgabe~3.3 --- Wechselkursänderung im Exportgeschäft (14 Punkte)}
\addcontentsline{toc}{subsection}{Aufgabe~3.3 --- Wechselkursänderung im Exportgeschäft}

Sie sind Mitarbeiter/-in im Vertrieb der \emph{Nordlicht Optik GmbH}, Kiel. Das
Unternehmen verkauft Laser-Messgeräte in die USA und nach Österreich. Ein Messgerät wird
zum Preis von 45.000~€ kalkuliert. Für die Fertigung bezieht Nordlicht Optik Sensoren von
einem Lieferanten in den USA; eine Jahreslieferung kostet 32.400~USD. Zu Jahresbeginn gilt der
Kurs 1~€ $=$ 1,08~USD (Mengennotierung).

\begin{enumerate}[label=\alph*)]
  \item \textbf{Berechnen} Sie den Preis eines Messgeräts in US-Dollar und den
    Euro-Betrag der Sensorlieferung zum Kurs von 1,08~USD. \hfill (2~P.)
  \item Zum Jahresende notiert der Kurs bei 1~€ $=$ 1,20~USD. \textbf{Stellen} Sie
    \textbf{fest}, ob der Euro auf- oder abgewertet hat. \textbf{Berechnen} Sie den neuen
    US-Dollar-Preis des Messgeräts (absolut und prozentuale Veränderung) sowie den neuen
    Euro-Betrag der Sensorlieferung. \hfill (5~P.)
  \item Ein wichtiger US-Kunde akzeptiert keine Preiserhöhung. Nordlicht Optik hält
    deshalb den bisherigen Dollarpreis aus a). \textbf{Berechnen} Sie den Euro-Erlös je
    Messgerät beim neuen Kurs und die Erlösminderung in Euro und Prozent. \hfill (2~P.)
  \item \textbf{Nennen} Sie zwei Möglichkeiten, wie sich Nordlicht Optik künftig gegen
    solche Wechselkursrisiken absichern kann. \hfill (2~P.)
  \item \textbf{Erläutern} Sie, warum bei den Lieferungen nach Österreich kein
    Wechselkursrisiko besteht, und \textbf{nennen} Sie zwei weitere Vorteile, die der
    Nordlicht Optik GmbH aus dem Europäischen Binnenmarkt bzw. der Währungsunion
    entstehen. \hfill (3~P.)
\end{enumerate}

\subsection*{Aufgabe~3.4 --- Regionaler Arbeitsmarkt (13 Punkte)}
\addcontentsline{toc}{subsection}{Aufgabe~3.4 --- Regionaler Arbeitsmarkt}

Sie sind Mitarbeiter/-in der \emph{Wirtschaftsförderungsgesellschaft Kreis Hollenbeck mbH}
und bereiten den jährlichen Arbeitsmarktbericht vor. Es liegen folgende Daten vor:

\begin{center}
\begin{tabular}{lrr}
\toprule
 & \textbf{Vorjahr} & \textbf{Berichtsjahr}\\
\midrule
Erwerbstätige & 236.800 & 234.000\\
registrierte Arbeitslose & 13.200 & 16.000\\
\bottomrule
\end{tabular}
\end{center}

Von den 16.000 Arbeitslosen des Berichtsjahres sind nach einer Auswertung der
Arbeitsagentur 7.400 Personen strukturell bedingt arbeitslos.

\begin{enumerate}[label=\alph*)]
  \item \textbf{Berechnen} Sie die Arbeitslosenquote für beide Jahre und die Veränderung
    in Prozentpunkten. \hfill (4~P.)
  \item \textbf{Ordnen} Sie die folgenden vier Fälle jeweils einer Form der
    Arbeitslosigkeit \textbf{zu} und \textbf{begründen} Sie kurz. \hfill (4~P.)
    \begin{enumerate}[label=(\arabic*)]
      \item Ein Maurer wird jedes Jahr von Dezember bis März nicht beschäftigt.
      \item Eine Industriemechanikerin eines Autozulieferers wird entlassen, weil die
        gesamtwirtschaftliche Nachfrage eingebrochen ist und die Aufträge fehlen.
      \item Ein Drucker verliert seine Stelle, weil seine Zeitungsdruckerei wegen der
        Digitalisierung dauerhaft schließt; sein Beruf wird in der Region kaum noch
        nachgefragt.
      \item Eine Bürokauffrau hat gekündigt, um in eine andere Stadt zu ziehen, und hat
        dort nach sechs Wochen Suche bereits eine Zusage für eine neue Stelle.
    \end{enumerate}
  \item \textbf{Berechnen} Sie den Anteil der strukturellen Arbeitslosigkeit an der
    gesamten Arbeitslosigkeit und \textbf{erläutern} Sie, warum ein hoher Anteil für den
    Landkreis besonders problematisch ist. \hfill (2~P.)
  \item \textbf{Schlagen} Sie je eine geeignete Maßnahme gegen die strukturelle und gegen
    die konjunkturelle Arbeitslosigkeit \textbf{vor} und \textbf{ordnen} Sie diese der
    aktiven bzw. passiven Arbeitsmarktpolitik \textbf{zu}. \hfill (3~P.)
\end{enumerate}

\subsection*{Aufgabe~3.5 --- Leitzinssenkung der EZB (12 Punkte)}
\addcontentsline{toc}{subsection}{Aufgabe~3.5 --- Leitzinssenkung der EZB}

Sie sind Mitarbeiter/-in in der Finanzierungsabteilung der \emph{Wohnbau Lindner GmbH},
eines Bauträgers für Eigentumswohnungen. Die Inflationsrate im Euroraum (HVPI) liegt seit
mehreren Monaten bei rund 1,3~\%; die Prognosen erwarten auch mittelfristig einen Wert
deutlich unter 2~\%. Die Wirtschaft im Euroraum stagniert. Der EZB-Rat beschließt, die drei
Leitzinsen um jeweils 0,25~Prozentpunkte zu senken.

\begin{enumerate}[label=\alph*)]
  \item \textbf{Erläutern} Sie, warum die EZB die Zinsen senkt, obwohl die Preise im
    Euroraum weiterhin steigen. \hfill (3~P.)
  \item \textbf{Stellen} Sie die Wirkungskette der Zinssenkung bis zur Auftragslage der
    Wohnbau Lindner GmbH \textbf{dar} (mindestens fünf Stufen). \hfill (5~P.)
  \item Ein Kollege meint: „Endlich hat die Bundesregierung die Zinsen gesenkt, um die
    Bauwirtschaft zu retten.“ \textbf{Nehmen} Sie zu dieser Aussage \textbf{Stellung}.
    \hfill (2~P.)
  \item \textbf{Erläutern} Sie zwei Gründe, warum die Zinssenkung möglicherweise nicht zu
    mehr Wohnungskäufen führt. \hfill (2~P.)
\end{enumerate}

\subsection*{Aufgabe~3.6 --- Konjunkturprogramm und Schuldenbremse (11 Punkte)}
\addcontentsline{toc}{subsection}{Aufgabe~3.6 --- Konjunkturprogramm und Schuldenbremse}

Sie sind Assistent/-in der Geschäftsführung der \emph{Brenner Tiefbau GmbH}. Deutschland
befindet sich in einer Rezession. In der Bundesregierung werden folgende Maßnahmen
diskutiert:

\begin{enumerate}[label=(\arabic*)]
  \item Bereits geplante Investitionen in Brücken und Schienenwege werden um zwei Jahre
    vorgezogen.
  \item Der Regelsteuersatz der Umsatzsteuer wird angehoben.
  \item Unternehmen erhalten befristet eine Sonderabschreibung auf neue Maschinen.
  \item Subventionen werden gekürzt; die Einsparungen werden einer Rücklage zugeführt.
\end{enumerate}

\begin{enumerate}[label=\alph*)]
  \item \textbf{Ordnen} Sie die vier Maßnahmen als expansiv oder restriktiv \textbf{zu}
    und \textbf{geben} Sie \textbf{an}, welche davon in der aktuellen Lage
    konjunkturgerecht sind. \hfill (4~P.)
  \item Das nominale BIP beträgt 4.400~Mrd.~€. Der Bund plant für das Konjunkturprogramm
    eine Nettokreditaufnahme von 41~Mrd.~€, davon sind 22~Mrd.~€ als konjunkturbedingt
    anerkannt (Konjunkturkomponente). Sonderregelungen (z.\,B. für Verteidigungsausgaben
    oder Sondervermögen) bleiben außer Betracht. \textbf{Prüfen} Sie rechnerisch, ob die
    Schuldenbremse des Grundgesetzes eingehalten wird. \hfill (4~P.)
  \item \textbf{Erläutern} Sie, warum das Investitionsprogramm (1) trotz guter Absicht
    prozyklisch wirken kann, und \textbf{erklären} Sie den Begriff Crowding-out.
    \hfill (3~P.)
\end{enumerate}

\subsection*{Aufgabe~3.7 --- Zahlungsbilanz und Handelskonflikt (14 Punkte)}
\addcontentsline{toc}{subsection}{Aufgabe~3.7 --- Zahlungsbilanz und Handelskonflikt}

Sie sind Mitarbeiter/-in der Exportabteilung der \emph{Weinmann Landtechnik GmbH}, die
Traktoren und Erntemaschinen weltweit verkauft. Für eine Marktanalyse liegen Ihnen
folgende (gerundete, angenommene) Zahlungsbilanzdaten eines Jahres für Deutschland vor
(in Mrd.~€):

\begin{center}
\begin{tabular}{lr}
\toprule
\textbf{Position} & \textbf{Mrd.~€}\\
\midrule
Warenexporte & 1.480\\
Warenimporte & 1.290\\
Einnahmen aus Dienstleistungen & 395\\
Ausgaben für Dienstleistungen & 455\\
Saldo der Primäreinkommen & $+115$\\
Saldo der Sekundäreinkommen & $-72$\\
\bottomrule
\end{tabular}
\end{center}

\begin{enumerate}[label=\alph*)]
  \item \textbf{Berechnen} Sie den Saldo der Handelsbilanz, der Dienstleistungsbilanz und
    der Leistungsbilanz. \hfill (4~P.)
  \item \textbf{Erläutern} Sie, in welcher Teilbilanz sich der Leistungsbilanzsaldo
    spiegelbildlich niederschlägt und warum die Zahlungsbilanz formal stets ausgeglichen
    ist. \hfill (2~P.)
  \item Der (fiktive) Staat \emph{Westland}, ein wichtiger Absatzmarkt der Weinmann
    Landtechnik GmbH, kündigt an: (1)~einen Einfuhrzoll von 20~\% auf Landmaschinen aus
    der EU, (2)~eine Höchstmenge von 500 importierten Traktoren pro Jahr, (3)~eine neue,
    aufwendige Sicherheitszertifizierung, die nur von Prüfstellen in Westland
    durchgeführt werden darf, (4)~Exportsubventionen für die eigenen
    Landmaschinenhersteller. \textbf{Unterscheiden} Sie tarifäre und nichttarifäre
    Handelshemmnisse und \textbf{ordnen} Sie die vier Maßnahmen \textbf{zu}. \hfill (4~P.)
  \item In Deutschland wird gefordert, die EU solle mit Gegenzöllen auf Waren aus Westland
    antworten. \textbf{Beurteilen} Sie diese Forderung, indem Sie je zwei Argumente dafür
    und dagegen anführen. \hfill (4~P.)
\end{enumerate}

\subsection*{Aufgabe~3.8 --- Emissionshandel im Unternehmen (11 Punkte)}
\addcontentsline{toc}{subsection}{Aufgabe~3.8 --- Emissionshandel im Unternehmen}

Sie sind Mitarbeiter/-in im Controlling der \emph{Hagedorn Glaswerke GmbH}. Das Werk
nimmt am Emissionshandel der EU teil. Für das kommende Jahr gelten folgende Planwerte:

\begin{itemize}
  \item kostenlos zugeteilte Emissionsberechtigungen: 40.000 (1 Berechtigung berechtigt
    zum Ausstoß von 1~t CO\textsubscript{2}),
  \item erwartete Emissionen ohne Gegenmaßnahmen: 46.000~t CO\textsubscript{2},
  \item angenommener Marktpreis: 80~€ je Berechtigung (konstant),
  \item Investitionsalternative: Eine Wärmerückgewinnungsanlage senkt die Emissionen um
    9.000~t CO\textsubscript{2} pro Jahr; sie verursacht jährliche Kosten (Abschreibung,
    Zinsen, Betrieb) von 540.000~€.
\end{itemize}

\begin{enumerate}[label=\alph*)]
  \item \textbf{Erläutern} Sie die Funktionsweise des Emissionshandels und \textbf{nennen}
    Sie das umweltpolitische Prinzip, auf dem er beruht. \hfill (3~P.)
  \item \textbf{Berechnen} Sie die jährliche Kostenbelastung aus dem Emissionshandel ohne
    und mit Wärmerückgewinnungsanlage und \textbf{entscheiden} Sie begründet, welche
    Alternative wirtschaftlich günstiger ist. \hfill (5~P.)
  \item \textbf{Beurteilen} Sie den Emissionshandel im Vergleich zu einer staatlichen
    Auflage (fester Emissionsgrenzwert je Werk), indem Sie zwei Vorteile und einen
    Nachteil des Emissionshandels darlegen. \hfill (3~P.)
\end{enumerate}
